\section*{Capítulo \ref{ch:b3:renal-osmoregulation} --- Fisiología renal y osmorregulación}

\begin{solution}{exo:b3:renal-osmoregulation:1}
\omterm{def:b3:renal-osmoregulation:nephron}{Glomérulo}: filtra el plasma menos las proteínas. \omterm{def:b3:renal-osmoregulation:nephron}{Túbulo proximal}:
reabsorbe dos tercios de la sal y del agua, toda la glucosa y los
aminoácidos, y el bicarbonato. Rama descendente de Henle: cede agua a
la médula. Rama ascendente: bombea NaCl hacia fuera, impermeable al
agua. Túbulo distal: ajusta el sodio (\omterm{def:b3:renal-osmoregulation:controls}{aldosterona}) y secreta potasio.
\omterm{def:b3:renal-osmoregulation:nephron}{Tubo colector}: reabsorción de agua bajo la vasopresina, manejo de
protones y de urea.
\end{solution}

\begin{solution}{exo:b3:renal-osmoregulation:2}
Aclaramiento: el volumen de plasma depurado de la sustancia por unidad
de tiempo, $UV/P$. Aquí, $100\times 1.5/2 = \qty{75}{mL/min}$, por
debajo de la TFG de \qty{120}{mL/min}: la sustancia se filtra y se
reabsorbe en parte (alrededor del \qty{38}{\%} de la cantidad
filtrada).
\end{solution}

\begin{solution}{exo:b3:renal-osmoregulation:3}
Un pez de agua dulce es más salado que su medio: el agua entra por las
branquias por ósmosis, de modo que beber solo añadiría a la inundación;
excreta una orina copiosa y diluida y las \omterm{def:b3:renal-osmoregulation:comparative}{células de cloruro} de sus
branquias captan Na$^{+}$ y Cl$^{-}$ activamente del agua. Un pez
marino es menos salado que el mar: pierde agua por las branquias, bebe
agua de mar para reponerla, y las \omterm{def:b3:renal-osmoregulation:comparative}{células de cloruro} expulsan la sal
tragada.
\end{solution}

\begin{solution}{exo:b3:renal-osmoregulation:4}
El \omterm{prop:b3:renal-osmoregulation:countercurrent}{efecto simple} es la diferencia de \qty{200}{mOsm/L} que la rama
ascendente gruesa mantiene entre el intersticio y su propio líquido al
bombear NaCl hacia fuera sin dejar que el agua lo siga. La rama
ascendente ha de ser impermeable al agua, pues de otro modo el agua
seguiría a la sal y la diferencia se desvanecería; la rama descendente
ha de ser permeable al agua para que el líquido que llega al recodo ya
haya sido concentrado por el intersticio y cada etapa trabaje sobre un
líquido más salado que la anterior. Dos ramas con la misma
permeabilidad no podrían multiplicar.
\end{solution}

\begin{solution}{exo:b3:renal-osmoregulation:5}
Extremo aferente: $60 - 18 - 28 = \qty{14}{mmHg}$; extremo eferente: $60 - 18 -
35 = \qty{7}{mmHg}$. A medida que se filtra agua, las proteínas que
quedan atrás se concentran y la presión oncótica sube a lo largo del
capilar; si alcanza \qty{42}{mmHg} la presión neta es cero y la
filtración se detiene antes del final (equilibrio de filtración). La
TFG depende entonces del flujo plasmático: un flujo más rápido
concentra las proteínas más despacio, de modo que la filtración
continúa a lo largo de un tramo más largo y la TFG sube con el flujo.
\end{solution}

\begin{solution}{exo:b3:renal-osmoregulation:6}
Carga filtrada: $\qty{40}{mL/min}\times\qty{0.14}{mmol/mL} =
\qty{5.6}{mmol/min}$. Excreción: $70\times 1 = \qty{0.07}{mmol/min}$.
Reabsorbido: $5.53/5.6 = \qty{98.75}{\%}$. Aclaramiento de sodio: $70\times
1/140 = \qty{0.5}{mL/min}$.
\end{solution}

\begin{solution}{exo:b3:renal-osmoregulation:7}
$900/1200 = \qty{0.75}{L}$ al día. A \qty{400}{mOsm/L}: $900/400 =
\qty{2.25}{L}$ de orina, de modo que hay que beber unos
$2.25 + 0.9 = \qty{3.15}{L}$ (menos el agua de los alimentos): el
paciente cuyo mecanismo de concentración falla ha de beber, y corre
peligro si no puede.
\end{solution}

\begin{solution}{exo:b3:renal-osmoregulation:8}
(a) Diabetes insípida: \qtyrange{10}{15}{L} de orina a
\qtyrange{50}{100}{mOsm/L}, sodio plasmático alto siempre que la bebida
se retrase, sed implacable. (b) Vasopresina inadecuada: orina escasa y
concentrada, agua retenida, sodio plasmático bajo por dilución, sed no
aumentada aunque se siga bebiendo ---confusión y convulsiones cuando el
sodio cae mucho. (c) Exceso de \omterm{def:b3:renal-osmoregulation:controls}{aldosterona}: sodio retenido con agua,
volumen expandido, presión arterial alta, potasio perdido (K$^{+}$
plasmático bajo), sodio plasmático casi normal, ya que el agua sigue a
la sal y el riñón escapa en parte; volumen de orina normal. (d) Dieta
sin sal: la renina y la \omterm{def:b3:renal-osmoregulation:controls}{aldosterona} suben en un día, el sodio urinario
cae a unos pocos milimoles al día, el sodio plasmático se mantiene, y
el volumen y la presión quedan algo más bajos.
\end{solution}

\begin{solution}{exo:b3:renal-osmoregulation:9}
Orina diluida: aclaramiento osmolar $100\times 6/300 = \qty{2}{mL/min}$;
\omterm{prop:b3:renal-osmoregulation:water}{aclaramiento de agua libre} $6 - 2 = +\qty{4}{mL/min}$: el riñón está
desechando agua sin soluto, como tras una sobrecarga de agua con la
vasopresina apagada. Orina concentrada: aclaramiento osmolar $1000\times 0.6/300 =
\qty{2}{mL/min}$; \omterm{prop:b3:renal-osmoregulation:water}{aclaramiento de agua libre} $0.6 - 2 = -\qty{1.4}{mL/min}$:
se devuelven al cuerpo \qty{1.4}{mL/min} de agua libre
---deshidratación, con la vasopresina alta.
\end{solution}

\begin{solution}{exo:b3:renal-osmoregulation:10}
Numérense las etapas desde la corteza; sea el intersticio de la etapa
$k$ igual a $I_{k}$. El líquido descendente en $k$ vale $I_{k}$; el líquido
ascendente en $k$ vale $I_{k} - 200$; y el líquido ascendente en $k$ es
el que salió del recodo y pasó por las etapas $n, n-1, \dots, k+1$. En
estado estacionario el líquido que entra a \qty{300}{} abandona lo alto
de la rama ascendente a $I_{1} - 200$, y el intersticio de cada etapa
supera al de la de encima como mucho en el \omterm{prop:b3:renal-osmoregulation:countercurrent}{efecto simple}, de modo que
$I_{k} \le 300 + 200k$, con igualdad cuando todas las etapas trabajan a
pleno efecto: el recodo alcanza $300 + 200n$. El gradiente está, por
tanto, acotado por la longitud del asa (el número de etapas) por el
\omterm{prop:b3:renal-osmoregulation:countercurrent}{efecto simple}, y no por la potencia de bombeo de célula alguna. $n$ es
la longitud del asa dividida por la distancia a lo largo de la cual el
líquido se equilibra con el intersticio: las asas yuxtamedulares
largas, y las larguísimas de los roedores del desierto, dan un $n$
grande; la capacidad de la bomba y el lavado del soluto por los vasos
rectos ponen tope al efecto.
\end{solution}

\begin{solution}{exo:b3:renal-osmoregulation:11}
\omterm{def:b3:renal-osmoregulation:comparative}{Agua metabólica}: $5\times 0.8 = \qty{4}{g}$ de almidón dan $4\times 0.6
= \qty{2.4}{g}$. Pérdidas: orina, $3/5500\ \unit{L} = \qty{0.55}{mL}$;
evaporación, \qty{1.5}{g} (ya neta de la recuperación de la nariz).
Pérdida total \qty{2.05}{g} frente a una ganancia de \qty{2.4}{g}: un
excedente de \qty{0.35}{g}, que cubre la pequeña pérdida fecal, y el
agua higroscópica de las semillas es un extra. Sobrevive sin beber. Sin
la recuperación nasal la pérdida respiratoria sería de
$1.5/0.3 = \qty{5}{g}$, y moriría en unos pocos días.
\end{solution}

\begin{solution}{exo:b3:renal-osmoregulation:12}
En el flujo paralelo la sangre y el líquido de diálisis se equilibran a
medio camino de la membrana y la diferencia de concentración que
impulsa la difusión cae a cero: en el mejor de los casos la sangre sale
con la mitad de su urea. En el flujo a contracorriente el líquido de
diálisis más fresco se encuentra con la sangre más depurada y el
gradiente se mantiene a lo largo de toda la membrana, de modo que la
sangre puede salir con casi la concentración de entrada del líquido de
diálisis ---el mismo principio que el \omterm{def:b3:renal-osmoregulation:nephron}{asa de Henle}. Aclaramiento: $300\times 0.7 =
\qty{210}{mL/min}$ durante $3\times 4 = \qty{12}{h}$ por semana, es
decir, $210\times 720 = \qty{151}{L}$ por semana; dos riñones a
\qty{125}{mL/min} depuran $125\times\num{10080} = \qty{1260}{L}$. La
máquina hace alrededor de la octava parte del trabajo de los riñones,
unos \qty{15}{mL/min} promediados en el tiempo ---bastante para vivir,
no bastante para estar bien.
\end{solution}

\begin{solution}{pb:b3:renal-osmoregulation:1}
\textbf{1.} Presión neta $55 - 15 - 30 = \qty{10}{mmHg}$; TFG $= 12.5
\times 10 = \qty{125}{mL/min}$.
\textbf{2.} $125\times 1440 = \qty{180}{L/d}$; $180/3 = 60$ veces.
\textbf{3.} Inulina: $62.5\times 1/0.5 = \qty{125}{mL/min}$, igual a la
TFG. Creatinina: $1.3\times 1/0.01 = \qty{130}{mL/min}$, algo mayor
porque el \omterm{def:b3:renal-osmoregulation:nephron}{túbulo proximal} secreta una pequeña cantidad de creatinina
además de la filtrada.
\textbf{4.} Aclaramiento de PAH, $13\times 1/0.02 = \qty{650}{mL/min}$,
el flujo plasmático renal; flujo sanguíneo,
$650/(1 - 0.45) = \qty{1180}{mL/min}$, unos \qty{1.2}{L/min}; fracción
de filtración, $125/650 = 0.19$.
\textbf{5.} Presión neta \qty{15}{mmHg}, TFG \qty{188}{mL/min} (menos en
realidad, ya que la presión oncótica sube a lo largo del capilar). La
angiotensina II contrae la arteriola eferente y eleva la presión
glomerular; bloquearla baja un poco la presión y la TFG. En la diabetes
los \omterm{def:b3:renal-osmoregulation:nephron}{glomérulos} filtran a presión alta y el filtro se desgasta; bajar la
presión retrasa el daño en años.
\textbf{6.} Los podocitos y sus diafragmas de hendidura (o la carga de
la membrana basal): la capa que retiene la albúmina. Perder
\qty{3}{g/d} baja la presión oncótica plasmática; el líquido sale de
los capilares en todas partes y los tejidos se hinchan ---el edema del
síndrome nefrótico.
\textbf{7.} Filtrada, $125\times 5 = \qty{0.625}{mmol/min}$; excretada,
$250\times 1 = \qty{0.25}{mmol/min}$; reabsorbida, $0.375$, es decir, el
\qty{60}{\%}; aclaramiento, $250/5 = \qty{50}{mL/min}$.
\textbf{8.} Filtrados $180\times 0.14 = \qty{25.2}{mol/d}$ de Na$^{+}$,
\qty{580}{g} de sodio, \qty{1.47}{kg} como NaCl. Excretados $100\times 1.44
= \qty{144}{mmol/d}$; reabsorbido $1 - 144/\num{25200} = \qty{99.4}{\%}$.
Se perdería casi kilo y medio de sal al día.
\textbf{9.} Filtrada, $125\times 5 = \qty{0.625}{mmol/min}$, es decir,
\qty{112}{mg/min}, \qty{162}{g/d}; no se excreta nada, ya que está por
debajo de $T_{m}$. La carga filtrada alcanza $T_{m}$ en $2.1/0.125 =
\qty{16.8}{mmol/L}$ (\qty{300}{mg/dL}) ---más alto que el umbral
observado, de unos \qty{10}{mmol/L}, porque las \omterm{def:b3:renal-osmoregulation:nephron}{nefronas} difieren y las
más débiles se derraman primero (la dispersión de la curva de
titulación).
\textbf{10.} Filtrada, $125\times 20 = \qty{2.5}{mmol/min}$; excretada,
$2.5 - 2.1 = \qty{0.4}{mmol/min}$, \qty{72}{mg/min}, $576$ mmol o
\qty{104}{g} al día. Orina extra, $576/300 = \qty{1.9}{L}$. La glucosa
arrastra agua a la orina (diuresis osmótica), el paciente se deshidrata
y tiene sed; \qty{104}{g} de glucosa son \qty{1.8}{MJ} perdidos al día,
y las células, incapaces de captar glucosa, queman grasa y proteína: el
peso cae.
\textbf{11.} Reabsorbidos $\approx \qty{25}{mol/d}$; ATP, $25/3 =
\qty{8.4}{mol/d}$; $8.4\times 50 = \qty{420}{kJ}$, alrededor del
\qty{6}{\%} de \qty{7}{MJ}.
\textbf{12.} Para secretar cada desecho, el túbulo necesitaría un
transportador que lo reconociese ---miles de ellos, y ninguno para una
molécula nunca antes encontrada. Filtrarlo todo lo pequeño y reabsorber
las pocas centenas de especies que el cuerpo quiere solo necesita
transportadores para lo que se sabe valioso; todo lo no reconocido se
marcha por omisión. El precio es la energía de la pregunta~11.
\textbf{13.} Aclaramiento osmolar, $600\times 1/290 = \qty{2.07}{mL/min}$;
\omterm{prop:b3:renal-osmoregulation:water}{aclaramiento de agua libre}, $1 - 2.07 = -\qty{1.07}{mL/min}$: el riñón
está conservando agua.
\textbf{14.} Mínimo, $600/1200 = \qty{0.5}{L}$; máximo, $600/50 =
\qty{12}{L}$; necesidad diaria en el mínimo, $0.5 + 0.9 = \qty{1.4}{L}$.
\textbf{15.} \qty{1000}{mOsm} a \qty{1200}{mOsm/L} necesitan
\qty{0.83}{L} de orina: ganancia neta de \qty{0.17}{L} antes del soluto
del día. Después: el soluto total del día es de \qty{1600}{mOsm},
\qty{1.33}{L} de orina, más \qty{0.9}{L} insensibles, \qty{2.23}{L} de
salida por \qty{1}{L} de entrada ---un déficit de \qty{1.23}{L}, frente
a \qty{1.4}{L} sin beber. El litro de agua de mar mejora el balance en
solo \qty{0.17}{L}.
\textbf{16.} Orina a \qty{50}{mOsm/L} para el día: $600/50 =
\qty{12}{L}$; hay que beber \qty{12.9}{L} cada día, unos \qty{13}{L}.
\textbf{17.} $B$ litros de cerveza aportan $20B$ mOsm; la orina máxima
es de $(600 + 20B)/50 = 12 + 0.4B$ litros, y debe arrastrar $B - 0.9$
litros: $B - 0.9 \le 12 + 0.4B$, $B \le \qty{21.5}{L}$. Unos veinte
litros ---pero si se come poco más, el soluto del día cae a
\qty{200}{mOsm} y el límite baja a unos \qty{8}{L}: la hiponatremia de
los bebedores empedernidos que no comen.
\textbf{18.} $298/290 - 1 = \qty{2.8}{\%}$: muy por encima del \qty{1}{\%}
que dispara la vasopresina y la sed; ambas son máximas. Con el soluto
constante: $290\times 42 = 298\times V$, $V = \qty{40.9}{L}$, un
déficit de alrededor de \qty{1.1}{L}.
\textbf{19.} Seis litros de agua diluyen el plasma mientras el sudor
retira sal; y el ejercicio, el dolor y el estrés liberan vasopresina
con independencia de la \omterm{def:b3:renal-osmoregulation:problem}{osmolaridad}, de modo que el riñón no puede
diluir la orina lo bastante para excretar el exceso. El sodio
plasmático cae, las células del encéfalo se hinchan ---ha habido
convulsiones y muertes. El riñón necesitaría la vasopresina apagada y
\qty{6}{L} de orina a $U_{\min}$; el corredor debería beber a demanda
de la sed y reponer la sal.
\textbf{20.} $3/5500\ \unit{L} = \qty{0.55}{mL}$ al día. Un riñón humano
necesitaría $3/1200 = \qty{2.5}{mL}$ para el mismo soluto: la rata
canguro emplea la quinta parte del agua.
\textbf{21.} Entradas: \qty{2.4}{g}. Salidas:
$1.6 + 0.3 + 0.55 = \qty{2.45}{g}$. Equilibrado con un margen de
\qty{0.05}{g} al día, que aporta el agua higroscópica de las semillas
(un pequeño porcentaje de \qty{5}{g}).
\textbf{22.} Agua, $500\times 0.65 = \qty{325}{L}$; la cuarta parte son
\qty{81}{L}; a \qty{5}{L/d}, dieciséis días.
\textbf{23.} $150/800 = \qty{0.19}{L}$, unos \qty{190}{mL} de salmuera
---menos que los \qty{100}{mL} de agua de mar y el agua del pescado
juntos, de modo que el ave gana agua. Su riñón, a \qty{600}{mOsm/L}
(\qty{300}{mmol/L} de NaCl), necesitaría \qty{0.5}{L} de orina para
arrastrar la misma sal, más agua de la que ingirió: un riñón que no
puede superar al agua de mar no permite a un ave beberla, y la
glándula, al doble de la concentración del mar, sí.
\textbf{24.} Entran \qty{30}{g} de agua, de modo que produce unos
\qty{30}{mL} de orina al día, a una pequeña fracción de la \omterm{def:b3:renal-osmoregulation:problem}{osmolaridad}
de la sangre (unas pocas decenas de mOsm/L frente a \num{300}). La sal
se va en esa orina y por difusión a través de las branquias; sin
captación activa por las \omterm{def:b3:renal-osmoregulation:comparative}{células de cloruro} desde un agua que contiene
un milimol por litro, el pez quedaría agotado en días.
\textbf{25.} TFG \qty{125}{mL/min}, \qty{180}{L} al día; orina
obligatoria \qty{0.5}{L} al día; un litro de agua de mar rinde unos
\qty{0.17}{L} netos una vez excretada su sal ---casi nada.
\end{solution}
